\subsection{The exact arithmetic-parameter classification for the two constants}

The rational-power rule in the preceding test can be
replaced by an arbitrary fixed nonzero algebraic
parameter and its extension powers. Retain the two
actual multiplication constants from VIII Figure~2,
p.~34 \cite{VIII}:
\[
 d(q)=q+2+q^{-1},\qquad
 c(q)=-q^2+q+q^{-1}-q^{-2}.
\]
No ambient transition coefficient is used in this
stronger scalar test.

Fix \(Q>1\), a prime power, a prime \(\ell\nmid Q\),
and a nonzero algebraic number \(\beta\).
Choose an \(\ell\)-adic embedding and a finite coefficient
field containing it. The arithmetic parameter rule is
\[
                              q_{\mathrm{src}}=\beta^n
\]
at extension degree \(n\). Thus one fixed parameter is
raised to extension powers; values are not chosen
independently at different degrees.

\begin{finding}[\fid{F65}, arbitrary fixed arithmetic parameter]
\label{f:jointArbitraryArithmeticParameter}
There exist bounded pure point complexes \(D,C\) over
\(\mathbb F_Q\), each of one integral Deligne weight
with possibly different weights, satisfying
\begin{equation}\label{eq:jointBetaEulerComparison}
 \operatorname{Tr}_{\mathrm{Eul}}(F^n;D)=d(\beta^n),
 \qquad
 \operatorname{Tr}_{\mathrm{Eul}}(F^n;C)=c(\beta^n)
\end{equation}
for every sufficiently large \(n\), if and only if
\(\beta^3=1\).

In those surviving cases \(c(\beta^n)\) vanishes
identically, and any pure complex \(C\) in
\eqref{eq:jointBetaEulerComparison} is zero.
The same classification holds if continuous arithmetic
Galois actions are required: the sufficient examples
have finite image.
\end{finding}

\begin{proof}
Finite power-sum uniqueness remains valid for eventual
equalities. For distinct nonzero roots
\(\alpha_1,\ldots,\alpha_s\), take \(s\) consecutive
sufficiently large degrees. The matrix of their powers
is a Vandermonde matrix times an invertible diagonal
matrix.

By BBD~5.1.8, pp.~126--127, a pure point complex of
weight \(w\) can contain a root of point weight \(m\)
only in degree \(m-w\) \cite{BBD}.
Its nonzero Euler multiplicity therefore has sign
\((-1)^{m-w}\). A fixed root cannot cancel between
different degrees. This allows Jordan blocks.
It also proves that a pure point complex with
identically zero Euler traces is zero: uniqueness
forces each grouped Euler multiplicity to vanish,
and its unique allowed degree then forces its
ordinary multiplicity to vanish.

Suppose first that \(\beta\ne1,-1\).
The three roots of
\[
                     d(\beta^n)=\beta^n+2+\beta^{-n}
\]
are distinct. Let \(w_D\) be the pure weight of \(D\).
The root \(1\) has weight zero and positive Euler
multiplicity \(2\), so its degree \(-w_D\) is even.
Thus \(w_D\) is even.
The root \(\beta\) has positive Euler multiplicity \(1\).
It must have an integral point weight \(m\), and its
degree \(m-w_D\) is even. Hence \(m\) is even.
The inverse root has weight \(-m\).

If the four numbers
\[
                    \beta^2,\quad\beta,\quad
                    \beta^{-1},\quad\beta^{-2}
\]
are distinct, their point weights \(2m,m,-m,-2m\)
are all even. But
\[
                 c(\beta^n)
                  =-\beta^{2n}+\beta^n+\beta^{-n}-\beta^{-2n}
\]
has coefficients \(-1,1,1,-1\). In any pure complex
of a fixed weight, the four allowed degrees have the
same parity. The mixed signs are impossible.

A collision among these four roots implies
\(\beta^r=1\) for some \(r\in\{1,2,3,4\}\):
these are the positive differences among the exponents
\(2,1,-1,-2\). We now check every case.

For \(\beta=1\), the sequences are \(d=4\) and \(c=0\).
For a primitive cube root,
\(\beta^2=\beta^{-1}\) and
\(\beta^{-2}=\beta\), so \(c\) again cancels identically.

For \(\beta=-1\),
\[
                         c(\beta^n)=2(-1)^n-2.
\]
The two distinct roots have weight zero and opposite
Euler signs, excluding purity.
For \(\beta\) of order four,
\[
                         c(\beta^n)
                         =\beta^n+\beta^{-n}-2(-1)^n.
\]
These three roots are distinct, all of weight zero,
and have both signs. Purity again fails.
This proves necessity and exhausts the possible
root collisions.

Conversely, if \(\beta^3=1\), choose \(D\) in
cohomological degree zero with eigenvalues
\(1,1,\beta,\beta^{-1}\), and choose \(C=0\).
It is pure of weight zero and satisfies the comparison
at every extension. The roots of unity give finite
image and hence a continuous arithmetic Galois action.
The earlier zero-trace argument proves the asserted
uniqueness of the zero object \(C\).
\end{proof}

\paragraph{The sufficient cases are genuinely degenerate.}
They even admit finite-etale scalar witnesses.
For \(\beta=1\), use four rational points.
For a primitive cube root, use
\[
              \operatorname{Spec}\mathbb F_Q
                  \ \sqcup\
              \operatorname{Spec}\mathbb F_{Q^3}.
\]
On geometric \(H^0\), Frobenius is one fixed point
and one three-cycle, with eigenvalues
\(1,1,\beta,\beta^{-1}\); higher cohomology vanishes.
This realizes \(d(\beta^n)\), while the zero object
realizes \(c(\beta^n)\).
It is a scalar witness, not a realization of the
source algebra, its canonical labels or a faithful
generic grading. Requiring a nonzero pure object
for the raw \(c\) coefficient eliminates even these
exceptions.

For \(\beta=Q^a\) with nonzero rational \(a\),
the parameter is not a root of unity, recovering
the preceding two-multiplication-constant obstruction.
The isolated \(\beta=Q^{1/2}\) pure realization of
\(c\) alone remains valid. The new finding asks for
the same parameter rule for \(d\) and \(c\); the
companion \(d\) then fails purity.

Separately retuned extension values, inserted-operator
traces, mixed coefficient complexes, and removal
of the raw vanishing factor are not covered.
VIII's signed/factored positivity is not being
replaced by an assertion that raw coefficients are
pure point traces.

% Separate paragraph for Part III.
\paragraph{The fixed arithmetic parameter in Grothendieck's language.}
A formal split-torus character and its assignment to
an arithmetic coefficient line are distinct data.
SGA~3, I, 4.7.2--3 and cote~90, scan~17 supply
action language; SGA~6, V, \S\S2--3 and cote~133,
scan~18 supply effective/virtual representation
language \cite{SGA3,SGA6,Archive}.
Neither supplies the fixed-parameter Euler comparison.

The bounded complexes fit SGA~6, I, Definition~4.7,
pp.~121--122 and Example~4.8, p.~122.
BBD~5.1.5 and~5.1.8, pp.~125--127 give the integral
point-weight and pure-complex degree conventions
\cite{SGA6,BBD}.
The Vandermonde argument, complete collision analysis
and finite-etale witnesses are proved here.
No \(\lambda\)-ring structure on the noncommutative
word algebra or unpublished purity assertion is used.

